% Source: 02 - Tue Sep 29 2026.pdf, page 13. Content remains in source order.
\sourcepage{02}{13}

\subsection{Search Problem }
The Search problem is a problem that asks whether a given element $m$ is contained in a given finite set $S\subseteq\mathbb N$.

It can be formalized as a concrete decision problem $L_{\mathrm{SEARCH}}$. 
\[
\begin{aligned}
\text{INSTANCE: }&\langle S,m\rangle:\ S\subseteq\mathbb N\text{ finite},\ m\in\mathbb N,\\
\text{QUESTION: }&m\in S?
\end{aligned}
\]

We can write the relative formal language  $L_{\mathrm{SEARCH}}$ as follows: 
\[
L_{\mathrm{SEARCH}}=\{x\in\bits^*:\ (x=\langle S,m\rangle)\land m\in S\}.
\]
Where x is a binary string, and x is in the language only if it is the encoding of a pair $\langle S,m\rangle$ where m is inside S. 

We can solve the search problem with a simple algorithm $A_{\mathrm{SEARCH}}$ that first checks if the input is a valid encoding, and then iterates through the elements of S to check if m is present. If m is found, it returns 1 (yes), otherwise it returns 0 (no).
\textbf{$A_{\mathrm{SEARCH}}(x)$}
\begin{tabbing}
\quad\=\quad\=\kill
\> if $(x\neq\langle S,m\rangle)$ then return $0$\\
\> $k\leftarrow|S|$\\
\> $*\ S=\{s_1,\ldots,s_k\}\ *$\\
\> for $i\leftarrow1$ to $k$ do\\
\>\> if $(s_i=m)$ then return $1$\\
\> return $0$
\end{tabbing}

Since we expect x to be a concise encoding, we can assume that the size of the set S is $|S|=O(|x|)$. Therefore, the algorithm runs in linear time with respect to the size of the input string x:
\[
T_{A_{\mathrm{SEARCH}}}(|x|)=O(|x|)\ \Rightarrow\
\{\mathrm{SEARCH},L_{\mathrm{SEARCH}}\}\in\classP.
\]
